
\chapter{Algèbres de Boole}
\begin{center}
    \includegraphics[width=4cm]{boole/img/boole.jpg}
\end{center}
\section{Définition d'une algèbre de Boole}

\begin{definition}[]
    Une \textit{algèbre de Boole}, c'est un ensemble $E$ muni de :
    \begin{itemize}
        \item 	deux lois \textit{binaires} notées + et $\times$;
        \item 	une loi \textit{unaire} qui à $a$ associe $\barmin{a}$;
        \item 	deux éléments particuliers notés 0 et 1;
    \end{itemize}
    et telle que , pour tous éléments a, b et c de E :
    \begin{itemize}
        \item 	Les deux lois binaires sont \textit{commutatives} :
              $a+b=b+a$ et $ab=ba$
        \item 	elles sont aussi associatives :
              $a+(b+c)=(a+b)+c$ et $a(bc)=(ab)c$
        \item 	$\times$ se distribue sur $+$ :
              $a(b+c)=ab+ac$
        \item 	+ se distribue sur $\times$ :
              $\boxed{a+bc=(a+b)(a+c)}$  \\[.5em]
              c'est la seule propriété contre-intuitive, on la notera $(*)$
        \item 	0 est neutre pour $+$: $0+a = a$
        \item 	0 est absorbant pour $\times$ : $0a=0$
        \item 	1 est neutre pour $\times$ : $1a=a$
        \item 	1 est absorbant pour $+$ : $1+a=1$
        \item 	on a $a+\barmin{a}=1$ et $a\barmin{a}=0$
    \end{itemize}
\end{definition}



\begin{exemple}[s]
    \begin{itemize}
        \item 	L'ensemble $\lbrace 0;1\rbrace$ muni de $\vee$, $\wedge$ et $\barmin{\ \ }$ est la plus simple des algèbres de Boole.
        \item 	Si on pose que deux propositions sont égales quand elles sont équivalentes, alors l'ensemble des propositions muni de $\vee$, $\wedge$ et $\barmin{\ \ }$  est une algèbre de Boole,
              0 étant la proposition fausse et 1 la proposition vraie.
        \item 	Comme on le verra au prochain chapitre, lorsqu'on se donne un ensemble $E$ et que l'on considère $\mathcal{P}(E)$, l'ensemble de ses parties, muni des opérations $\cap$ (intersection, joue le rôle de $\times$ ) et $\cup$ (union, joue le rôle de +), alors $P(E)$ est une algèbre de Boole, 0 étant l'\textit{ensemble vide} et 1 étant E lui-même.
    \end{itemize}
\end{exemple}

\begin{exercice}[]
    Montrer que :
    \begin{enumerate}
        \item 	$(a+b)(\barmin{a}+c)(\barmin{b}+\barmin{c})=a\barmin{b}c+\barmin{a}b\barmin{c}$
        \item 	$\barmin{a}\barmin{b}\barmin{c}\barmin{d}+a\barmin{b}\barmin{c}\barmin{d}+	\barmin{a}b\barmin{c}\barmin{d}+ab\barmin{c}\barmin{d}=\barmin{b}\barmin{c}$
        \item 	$ab+\barmin{a}\barmin{b}+ \barmin{a}b=\barmin{a}+b$
    \end{enumerate}
\end{exercice}

\begin{propriete}[s]
    \begin{itemize}
        \item 			Dans une algèbre de Boole, pas besoin de multiplication par un entier car pour tout élément $a$ :
              $$a+a = a$$
              Donc $2a = a$ et par suite pour tout entier naturel non nul $n$ : $na=a$.
        \item Pas besoin de puissances non plus  car pour tout élément $a$ :
              $$aa = a$$
              Donc $a^2 = a$ et par suite pour tout entier naturel non nul $n$ : $a^n=a$.
    \end{itemize}
\end{propriete}

\textbf{Preuve :}\\
Soit $a$ un élément de E:
\begin{multicols}{2}
    \begin{itemize}
        \item 	\begin{tabbing}
                  $a$ \= 	$=a+0$\\[.5em]
                  \>	$=a+a\barmin{a}\qquad$ et on utilise $(*)$\\[.5em]
                  \>	$=(a+a)(a+\barmin{a})$\\[.5em]
                  \>	$=(a+a)1$\\[.5em]
                  \>	$=a+a$
              \end{tabbing}
        \item 	\begin{tabbing}
                  $a$ \= $=a\times 1$\\[.5em]
                  \> $=a(a+\barmin{a})$\\[.5em]
                  \> $=aa+a\barmin{a}$\\[.5em]
                  \> $=aa+0$\\[.5em]
                  \> $=aa$
              \end{tabbing}
    \end{itemize}
\end{multicols}
\begin{propriete}
    Si 2 éléments $x$ et $y$ de E sont tels que $xy=0$ et $x+y=1$, alors $x=\barmin{y}$.
\end{propriete}
\textbf{Preuve :}
\begin{tabbing}
    $x+y=1\quad$	\=	$\Rightarrow\quad (x+y)\barmin{y}=\barmin{y}$\\[.5em]
    \>	$\Rightarrow\quad x\barmin{y}+y\barmin{y}=\barmin{y}$\\[.5em]
    \>	$\Rightarrow\quad x\barmin{y}+0=\barmin{y}$\\[.5em]
    \>	$\Rightarrow\quad x\barmin{y}+xy=\barmin{y}$\\[.5em]
    \>	$\Rightarrow\quad x(\barmin{y}+y)=\barmin{y}$\\[.5em]
    \>	$\Rightarrow\quad x\times 1=\barmin{y}$\\[.5em]
    \>	$\Rightarrow\quad x=\barmin{y}$
\end{tabbing}

\begin{propriete}[ : lois de De Morgan]
    Quels que soient les éléments de $x$ et $y$ de E on a $$\barmin{x\times y}=\barmin{x}+\barmin{y}  \qquad\textrm{et}\qquad \barmin{x+y}=\barmin{x}\times\barmin{y}$$
\end{propriete}

\textbf{Preuve de la première égalité}.\\
Si on pose $A=xy$ et $B=\barmin{x}+\barmin{y}$, alors on doit montrer que $\barmin{A}=B$.\\
On utilise la propriété précédente et on montre donc que $A+B = 1$ et que $AB=0$ :
\begin{multicols}{2}

    \begin{tabbing}
        $B+A$ 	\=	$=\barmin{x}+\barmin{y}+xy$ et on utilise $(*)$\\[.5em]
        \>	$=(\barmin{x}+\barmin{y}+x)(\barmin{x}+\barmin{y}+y)$\\[.5em]
        \>  $=(1+\barmin{y})(1+\barmin{x})$\\[.5em]
        \>  $=1\times 1$\\[.5em]
        \> 	$=1$
    \end{tabbing}

    \begin{tabbing}
        $AB$ \=  $=(\barmin{x}+\barmin{y})xy$\\[.5em]
        \>  $=\barmin{x}xy+\barmin{y}xy$\\[.5em]
        \> 	$=0y+0x$\\[.5em]
        \> 	$=0+0$\\[.5em]
        \> 	$=0$
    \end{tabbing}
\end{multicols}
\begin{exercice}[]
    Montrer la deuxième loi en faisant la même chose.
\end{exercice}



\section{Diagrammes de Karnaugh}

Cette méthode a été développée par Maurice Karnaugh en 1953. Elle fait appel à des tableaux qui représentent des expressions booléennes.\\
On colorie les cases qui correspondent aux différents termes de l'expression. À la fin on peut reconnaître et lire une expression simplifiée.

\begin{propriete}[]
    Un produit correspond à une \textit{intersection}, une somme à une \textit{union}.
\end{propriete}

\subsection{Avec 2 variables}

Le diagramme est un carré dont les quatre cases correspondent aux quatre produits $ab$, $a\barmin{b}$, $\barmin{a}b$, $\barmin{a}\barmin{b}$.

\begin{center}
    \begin{tabular}{|c|c|}
        \hline
        $ab$          & $a\barmin{b}$          \\
        \hline
        $\barmin{a}b$ & $\barmin{a}\barmin{b}$ \\
        \hline
    \end{tabular}
\end{center}
Ci dessous on représente a, b, $\barmin{a}$, $\barmin{b}$ et une dernière expression.
\begin{center}
    \includegraphics[width=\linewidth]{boole/img/karnaugh1.png}
\end{center}

Cette dernière peut se lire $a+\barmin{b}$ ou bien $ab+a\barmin{b}+\barmin{a}\barmin{b}$.

\begin{exercice}[]
    Propose une troisième interprétation du dernier diagramme.
\end{exercice}

\subsection{Avec 3 variables}

Le diagramme est un rectangle dont les huit cases correspondent aux produits suivants.

\begin{center}
    \begin{tabular}{|c|c|c|c|}

        \hline
        $abc$          & $ab\barmin{c}$          & $a\barmin{b}\barmin{c}$          & $a\barmin{b}c$          \\
        \hline
        $\barmin{a}bc$ & $\barmin{a}b\barmin{c}$ & $\barmin{a}\barmin{b}\barmin{c}$ & $\barmin{a}\barmin{b}c$ \\
        \hline
    \end{tabular}
\end{center}
\begin{center}
    \includegraphics[width=\linewidth]{boole/img/karnaugh2.png}\\

    \includegraphics[width=\linewidth]{boole/img/karnaugh3.png}
\end{center}

\begin{exercice}[]
    Faire les tables de Karnaugh des expressions $ab$, $a\barmin{b}$, $ac$, $a\barmin{c}$,  $\barmin{a}b$, $\barmin{a}\barmin{b}$, $\barmin{a}c$, $\barmin{a}\barmin{c}$, $bc$, $\barmin{b}c$, $b\barmin{c}$, $\barmin{b}\barmin{c}$.
\end{exercice}

\section{Exercices}

\begin{exercice}
    $\mathcal{B}$ est une algèbre de Boole.
    \begin{enumerate}
        \item 	Montrer par le calcul que $\forall a\in\mathcal{B},\forall b\in\mathcal{B},\; a+ab = a$;
        \item 	Montrer par le calcul que $\forall a\in\mathcal{B},\forall b\in\mathcal{B},\; a(a+b) = a$.
        \item 	Montrer par le calcul que $\forall a\in\mathcal{B},\forall b\in\mathcal{B},\; a+\barmin{a}b = a+b$.
    \end{enumerate}
    Vérifier ces trois égalités avec des tables de Karnaugh.
\end{exercice}

\begin{exercice}
    \'Ecrire de deux façons possibles l'expression booléenne représentée par le tableau de Karnaugh suivant.
    \begin{center}
        \includegraphics[width=4cm]{boole/img/ex29.png}
    \end{center}
\end{exercice}

\begin{exercice}
    \'Ecrire l'expression booléenne représentée par le tableau de Karnaugh suivant sous la forme d'une somme de deux variables booléennes prises parmi $x$, $y$, $z$, $\barmin{x}
    $, $\barmin{y}$ et $\barmin{z}$.
    \begin{center}
        \includegraphics[width=4cm]{boole/img/ex30.png}
    \end{center}
\end{exercice}

\begin{exercice}
    $x$,$y$ et $z$ sont trois élément d'une algèbre booléenne B.\\

    \'Ecrire l'expression $y\overline{(xy+z)}$ sous la forme d'un produit de trois variables booléennes prises parmi $x$, $y$, $z$, $\barmin{x}
    $, $\barmin{y}$ et $\barmin{z}$.\\
\end{exercice}


\begin{exercice}
    $a$,$b$ et $c$ sont trois élément d'une algèbre booléenne B.

    \begin{enumerate}
        \item 	\'Ecrire  l'expression $(a+\barmin{b}c)(b+\barmin{c})$ sous la forme d'une somme de deux produits de deux variables booléennes prises parmi $a$, $b$, $c$, $\barmin{a}
              $, $\barmin{b}$ et $\barmin{c}$.
        \item 	Représenter ce résultat dans une table de Karnaugh et en déduire une nouvelle expression.
    \end{enumerate}
\end{exercice}

\begin{exercice}
    B est une algèbre booléenne. On définit l'opération « nor »{}, notée $\downarrow$ par : $$\forall a\in\mathcal{B},\forall b\in\mathcal{B},\; a\downarrow b = \overline{a+b}$$

    Cette opération est dite \textit{universelle} car elle permet de retrouver toutes les autres opérations.

    \begin{enumerate}
        \item 	Montrer que $\forall a\in\mathcal{B}\; a\downarrow a = \overline{a}$.
        \item 	En déduire que $$\forall a\in\mathcal{B},\forall b\in\mathcal{B},\; (a\downarrow b)\downarrow(a\downarrow b) = a+b$$
        \item 	Comment à partir de $a$, $b$ et $\downarrow$ obtenir $ab$ ?\\
    \end{enumerate}
\end{exercice}

\begin{exercice}[ : Polynésie juin 2018]
    Une société de fabrication et d'installation de fibre optique a besoin de recruter un informaticien, femme ou homme. La direction des ressources humaines considère qu'une candidature est recevable lorsqu'elle satisfait à l'une au moins des conditions suivantes:

    \begin{itemize}
        \item le candidat est âgé de 25 ans ou moins et est titulaire du BTS SIO;
        \item le candidat est âgé de 25 ans ou moins, n'est pas titulaire du BTS SIO et possède de l'expérience;
        \item le candidat est âgé de strictement plus de 25 ans  et est titulaire du BTS SIO;
    \end{itemize}

    \smallskip
    On définit les variables booléennes $a$, $b$, $c$ de la façon suivante:
    \begin{itemize}
        \item $a=1$ si le candidat est âgé de strictement plus de 25 ans, $a=0$ sinon;
        \item $b=1$ si le candidat est titulaire d'un BTS SIO, $b=0$ sinon;
        \item $c=1$ si le candidat a de l'expérience, $c=0$ sinon.
    \end{itemize}

    \begin{enumerate}
        \item Écrire une expression booléenne $E$ traduisant qu'une candidature est recevable, à l'aide des variables booléennes $a$, $b$, $c$.
        \item À l'aide d'un tableau de Karnaugh, déterminer une écriture simplifiée de $E$ sous la forme d'une somme de deux termes. En déduire une interprétation simplifiée des conditions pour qu'une candidature soit recevable.
        \item Une candidate a 21 ans, aucune expérience, mais est titulaire du BTS SIO. Remplit-elle les critères de recrutement?
        \item Donner une expression simple de $\overline{E}$.\\
    \end{enumerate}
\end{exercice}

\begin{exercice}[ : Polynésie mai 2017]

    Cinq joueurs, notés A, B, C, D et E, jouent régulièrement à un jeu en ligne.

    Chaque partie de ce jeu oppose deux adversaires.

    \begin{enumerate}
        \item Dans cette question, on note $J = \lbrace A, B, C, D, E\rbrace$ l'ensemble des cinq joueurs.

              On note $V(x~;~y)$ le prédicat : « le joueur $x$ a déjà battu le joueur $y$ ».

              Ainsi, la valeur $V$(A ; B) est VRAI, et la valeur de $V$(B ; A) est FAUX.

              \smallskip

              On définit trois prédicats :


              \textbf{P1} : $\forall x \in J,\: \exists\, y \in J,\: x \neq y$  et $V(x~;~y)$

              \textbf{P2}: $\exists\, x \in J,\: \forall y \in J,\: x \neq y$  et $V(x~;~y)$

              \textbf{P3}: $\exists\, y \in J,\: \forall x \in J,\: x \neq y$  et $V(x~;~y)$

              \smallskip

              Associer à chaque prédicat P1, P2, P3, celle des trois phrases suivantes qui lui correspond parmi
              les phrases suivantes. Aucune justification n'est demandée.

              « Il existe un joueur qui a été battu par tous les autres joueurs ».

              « Tous les joueurs ont battu au moins un autre joueur  ».

              « Il existe un joueur qui a battu tous les autres joueurs  ».

        \item Un joueur reçoit un bonus lorsqu'il vérifie l'un au moins des trois critères suivants:

              $\bullet~~$ le joueur a participé à 20 parties ou davantage, et il a affronté plusieurs adversaires
              différents ;

              $\bullet~~$  le joueur n'a pas affronté plusieurs adversaires différents, et il a obtenu strictement plus de
              victoires que de défaites ;

              $\bullet~~$  le joueur n'a pas obtenu strictement plus de victoires que de défaites, et il a participé à 20
              parties ou davantage.

              On définit les variables booléennes $a$, $b$, $c$ de la façon suivante:

              $a = 1$ si le joueur a participé à 20 parties ou davantage ; $a = 0$ sinon;

              $b = 1$ si le joueur a affronté plusieurs adversaires différents ; $b = 0$ sinon;

              $c = 1$ si le joueur a obtenu strictement plus de victoires que de défaites ; $c = 0$ sinon.
              \begin{enumalph}
                  \item Écrire une expression booléenne $F$ traduisant les conditions permettant à un joueur d'obtenir le bonus.
                  \item À l'aide d'un tableau de Karnaugh ou d'un calcul booléen, déterminer une écriture simplifiée de $F$ sous forme d'une somme de deux termes.
                  \item En déduire une formulation simplifiée des critères permettant à un joueur d'obtenir le bonus.
              \end{enumalph}
    \end{enumerate}
\end{exercice}


\begin{exercice}[]
    Soient $a$, $b$ et $c$ trois éléments d'une algèbre booléenne $\mathcal{B}$.

    \begin{enumerate}
        \item 	Soient $A=ab+\barmin{c}$ et $B=\barmin{a}+bc$, montrer par le calcul que $$\barmaj{A}=\barmin{a}c+\barmin{b}c$$ et que $$\barmaj{B}=a\barmin{c}+a\barmin{b}$$
        \item 	Soit $C=\barmin{a}\,\barmin{b}\,\barmin{c}+a\barmin{b}c+\barmin{a}\,\barmin{b}+a\barmin{b}\,\barmin{c}$.\\
              Utiliser un diagramme de Karnaugh pour simplifier C.\\
    \end{enumerate}
\end{exercice}

\begin{exercice}[]
    Une salle dédiée à l'informatique va être aménagée au lycée.\\
    Le réseau qui équipera cette salle doit satisfaire au moins l'une des conditions suivantes :

    \begin{itemize}
        \item le réseau compte 5 postes ou plus et il existe un poste qui ne reçoit pas de données en entrée
        \item il existe un poste qui ne reçoit pas de données en entrée, et le réseau compte strictement moins  de 5 postes, et il comporte strictement plus de 12 connexions ;
        \item le réseau comporte 12 connexions ou moins.
    \end{itemize}

    On définit les variables booléennes suivantes:


    \begin{itemize}
        \item $a = 1$ si le réseau compte 5 postes ou plus, $a = 0$ sinon;
        \item $b = 1$ s'il existe un poste qui ne reçoit pas de données en entrée, $b = 0$ sinon;
        \item $c = 1$ si le réseau comporte 12 connexions ou moins, $c = 0$ sinon.
    \end{itemize}

    \begin{enumerate}
        \item Cette question est une question à choix multiple. Une seule réponse est correcte. Recopier sur la copie seulement la réponse correcte. On ne demande pas de justification.

              Parmi les quatre phrases suivantes, donner celle qui traduit la variable $\overline{b}$ :


              \begin{itemize}
                  \item réponse A : « il existe un poste qui reçoit des données en entrée  » ;
                  \item réponse B : « tout poste reçoit des données en entrée  » ;
                  \item réponse C : « il existe un poste qui envoie des données en sortie  » ;
                  \item réponse D : « aucun poste ne reçoit des données en entrée  ».
              \end{itemize}

        \item Donner l'expression booléenne E traduisant les critères voulus pour un réseau informatique.
        \item À l'aide d'un tableau de Karnaugh ou par des calculs, exprimer $E$ comme somme de deux
              variables booléennes.
        \item Traduire les critères de sélection simplifiés, à partir de l'expression obtenue à la question 3.
        \item Un réseau dans lequel 2 postes ne reçoivent pas de données en entrée et qui comporte 15
              connexions répond-il aux critères voulus ? Justifier la réponse.
    \end{enumerate}
\end{exercice}

\begin{exercice}[ : multiplexeur]
    Un multiplexeur est un peu comme une télécommande à 2 boutons : on peut le résumer par une expression booléenne $E$ dépendant de 3 variables $a$, $b$ et $c$.
    \begin{itemize}
        \item Lorsque $a=0$, la valeur de $E$ est celle de $b$.
        \item Lorsque $b=0$, la valeur de $E$ est celle de $c$.
    \end{itemize}
    \begin{enumerate}
        \item À l'aide de cette information, complète la table de vérité suivante.

              \begin{center}
                  \tabstyled
                  \begin{tabular}{|c|c|c|c|}
                      \hline
                      
                      \ccell a & \ccell b & \ccell c & \ccell E \\
                      \hline
                      0      & 0      & 0      &        \\
                      \hline
                      0      & 0      & 1      &        \\
                      \hline
                      0      & 1      & 0      &        \\
                      \hline
                      0      & 1      & 1      &        \\
                      \hline
                      1      & 0      & 0      &        \\
                      \hline
                      1      & 0      & 1      &        \\
                      \hline
                      1      & 1      & 0      &        \\
                      \hline
                      1      & 1      & 1      &        \\
                      \hline
                  \end{tabular}
              \end{center}
        \item Renseigne ensuite la table de Karnaugh
              \begin{center}
                  \begin{tikzpicture}
                      \draw (0,0) rectangle (1,1) rectangle(2,0) rectangle(3,1) rectangle(4,0);
                      \draw (0,0) rectangle (1,-1) rectangle(2,0) rectangle(3,-1) rectangle(4,0);
                  \end{tikzpicture}
              \end{center}
        \item Déduis-en une expression simple de $E$.

        \item Dessine le circuit électronique qui représente E.
    \end{enumerate}
\end{exercice}